% Options for packages loaded elsewhere
\PassOptionsToPackage{unicode}{hyperref}
\PassOptionsToPackage{hyphens}{url}
\PassOptionsToPackage{dvipsnames,svgnames,x11names}{xcolor}
%
\documentclass[
  10pt,
  letterpaper,
]{article}
\usepackage{amsmath,amssymb}
\usepackage{iftex}
\ifPDFTeX
  \usepackage[T1]{fontenc}
  \usepackage[utf8]{inputenc}
  \usepackage{textcomp} % provide euro and other symbols
\else % if luatex or xetex
  \usepackage{unicode-math} % this also loads fontspec
  \defaultfontfeatures{Scale=MatchLowercase}
  \defaultfontfeatures[\rmfamily]{Ligatures=TeX,Scale=1}
\fi
\usepackage{lmodern}
\ifPDFTeX\else
  % xetex/luatex font selection
\fi
% Use upquote if available, for straight quotes in verbatim environments
\IfFileExists{upquote.sty}{\usepackage{upquote}}{}
\IfFileExists{microtype.sty}{% use microtype if available
  \usepackage[]{microtype}
  \UseMicrotypeSet[protrusion]{basicmath} % disable protrusion for tt fonts
}{}
\makeatletter
\@ifundefined{KOMAClassName}{% if non-KOMA class
  \IfFileExists{parskip.sty}{%
    \usepackage{parskip}
  }{% else
    \setlength{\parindent}{0pt}
    \setlength{\parskip}{6pt plus 2pt minus 1pt}}
}{% if KOMA class
  \KOMAoptions{parskip=half}}
\makeatother
\usepackage{xcolor}
\usepackage[margin=0.78in]{geometry}
\usepackage{longtable,booktabs,array}
\usepackage{calc} % for calculating minipage widths
% Correct order of tables after \paragraph or \subparagraph
\usepackage{etoolbox}
\makeatletter
\patchcmd\longtable{\par}{\if@noskipsec\mbox{}\fi\par}{}{}
\makeatother
% Allow footnotes in longtable head/foot
\IfFileExists{footnotehyper.sty}{\usepackage{footnotehyper}}{\usepackage{footnote}}
\makesavenoteenv{longtable}
\setlength{\emergencystretch}{3em} % prevent overfull lines
\providecommand{\tightlist}{%
  \setlength{\itemsep}{0pt}\setlength{\parskip}{0pt}}
\setcounter{secnumdepth}{5}
\newcommand{\Z}{\mathbb Z}
\newcommand{\C}{\mathbb C}
\newcommand{\Pone}{\mathbb P^1}
\newcommand{\cO}{\mathcal O}
\newcommand{\Pic}{\operatorname{Pic}}
\newcommand{\NS}{\operatorname{NS}}
\newcommand{\rk}{\operatorname{rk}}
\newcommand{\im}{\operatorname{im}}
\newcommand{\coker}{\operatorname{coker}}
\newcommand{\Aut}{\operatorname{Aut}}
\newcommand{\Spec}{\operatorname{Spec}}
\newcommand{\dP}{\mathrm{dP}}
\ifLuaTeX
  \usepackage{selnolig}  % disable illegal ligatures
\fi
\usepackage{bookmark}
\IfFileExists{xurl.sty}{\usepackage{xurl}}{} % add URL line breaks if available
\urlstyle{same}
\hypersetup{
  pdftitle={Corrective Addendum and Mathematical Results Map},
  pdfauthor={Independent technical review},
  colorlinks=true,
  linkcolor={blue},
  filecolor={Maroon},
  citecolor={Blue},
  urlcolor={blue},
  pdfcreator={LaTeX via pandoc}}

\title{Corrective Addendum and Mathematical Results Map}
\usepackage{etoolbox}
\makeatletter
\providecommand{\subtitle}[1]{% add subtitle to \maketitle
  \apptocmd{\@title}{\par {\large #1 \par}}{}{}
}
\makeatother
\subtitle{Supplement to the full mathematical peer review of the
108-page S6 manuscript}
\author{Independent technical review}
\date{25 August 2026}

\begin{document}
\maketitle

\section{Corrective addendum dated 25 August
2026}\label{corrective-addendum-dated-25-august-2026}

This edition integrates the mathematical claims contained in the
user-supplied Codex progress log, but it does not treat that log as
evidence by itself. The log reports two strengthened proof lines: a
direct period-theoretic proof that

\[
a(X)=1,
\qquad
\mathcal M(X)=f^*\mathbb C(t),
\]

and an integral nearby-cycle/Leray computation that proves the
homology-sphere conclusion without the global cusp-collapse theorem. It
also reports corrections to a modular-frame factor, a Neron--Severi
simplification, Lemma A.3's variance, and the distinction between
\(\Lambda\) and \(\bar\Lambda\). Only the progress log, not the separate
Codex workbench files or its claimed 38-check harness, was supplied.
Every point promoted in this report was therefore re-located and
re-checked in the audited 108-page manuscript itself.

The result is a genuine but limited correction to the review of 24
August 2026.

\begin{enumerate}
\def\labelenumi{\arabic{enumi}.}
\tightlist
\item
  The direct algebraic-dimension argument in Section 9.1 is stronger
  than the earlier review stated. Its period calculation, very-general
  Neron--Severi computation, exclusion of algebraic dimensions one and
  two for the fibres, and meromorphic descent form a complete argument
  conditional only on the construction of the period family and standard
  structure theory for complex tori. It is independent of Section 7 and
  of Campana--Demailly--Peternell.
\item
  The second topology route in Section 7.7 and Appendix B is not merely
  corroborative. It is a logically sufficient topology engine for the
  recognition of \(X\) as \(S^6\), provided the based fundamental-group
  relations, the finite-fibre specialization lattices, the integral cusp
  specialization theorem, and the geometric identification of the Leray
  transgressions are correct. Proposition 7.2's global collapse is
  therefore not theorem-critical for the sphere-recognition chain,
  although it remains necessary for the manuscript's first topology
  route and for its explicit collapse-map and generator claims.
\item
  Lemma A.3 contains a confirmed variance error. The displayed
  Mayer--Vietoris sequence and every map subsequently computed are
  cohomological, although the proof writes \(H_k\). Replacing \(H_k\) by
  \(H^k\) throughout that calculation and applying universal
  coefficients at the end repairs the lemma without changing its stated
  groups.
\item
  The corrected source already contains the modular factor
  \(\mathsf e=F^{-1}e\) needed on source page 71, and the current
  algebra in Lemma 9.3 is correct. Those are repairs present in the
  audited text, not defects of the current version.
\item
  The current cusp notation correctly uses
  \(\bar\Lambda=\Lambda/\Lambda_{\mathrm{tor}}\) as the deck group. The
  Codex log does not identify the earlier erroneous occurrence, so no
  additional current-source error is asserted on that point.
\end{enumerate}

The recommendation remains \textbf{major revision and independent
specialist replication}. The reason is no longer a suspicion that the
global collapse is the sole route to the homology calculation, nor an
unresolved algebraic-dimension argument. The remaining burden is
concentrated in the actual geometric realization of the displayed local
and global maps, integral nearby-cycle functoriality and orientation
data, the based van Kampen and transgression comparisons, and the
analytic duality/base-change/literature-comparison chain in Section 10.

\subsection{Mathematical results so far: a cold-entry
guide}\label{mathematical-results-so-far-a-cold-entry-guide}

\subsubsection{The construction being
claimed}\label{the-construction-being-claimed}

The manuscript starts with a rank-four integral lattice \(\Lambda\) and
monodromies of orbifold orders \(3\) and \(4\), together with a
unipotent cusp monodromy. Period functions \(\tau,\mu,\beta\) define a
holomorphic family of two-dimensional complex tori over the
thrice-punctured projective line. Three local fillings are then
attached: a toric, reduced, irreducible, nonnormal normal-crossing fibre
\(W\) at the cusp, and multiple fibres \(3S_1\) and \(4S_2\) whose
reductions are bielliptic surfaces. The gluing parameters enter the
topology only through

\[
p=12\ell_0-4\ell_1-3\ell_2.
\]

For the actual data \((\ell_0,\ell_1,\ell_2)=(0,1,-1)\), one has
\(p=-1\).

\subsubsection{The topology chain now
isolated}\label{the-topology-chain-now-isolated}

The based van Kampen calculation produces

\[
\pi_1(X)=\langle c,x,y\mid c\ \text{central},\ xy=c^{\ell_0},\ x^3=c^{\ell_1},\ y^4=c^{\ell_2}\rangle
\cong \mathbb Z/|p|,
\]

conditional on the geometric identification of the three meridians and
their drift terms. Thus the actual value \(p=-1\) gives \(\pi_1(X)=1\).

The independent Leray route uses the exact cusp-specialization statement

\[
H^q(W;\mathbb Z)\xrightarrow{\ \operatorname{sp}^q\ }
H^q(F;\mathbb Z)^{T_0}
\]

as an isomorphism for \(0\le q\le4\), together with the finite-fibre
specialization indices. It obtains

\[
H^0(B,R^1f_*\mathbb Z)=12\mathbb Z\gamma,
\quad
H^0(B,R^2f_*\mathbb Z)=\mathbb Z\cdot2q,
\quad
H^0(B,R^3f_*\mathbb Z)=\mathbb Z\cdot2\gamma uw.
\]

The three relevant transgressions are

\[
\begin{aligned}
d_2^{0,1}(12\gamma)&=\pm p\,\omega,\\
d_2^{0,2}(2q)&=\pm p\,\omega\bar\delta,\\
d_2^{0,3}(2\gamma uw)&=\pm p\,\omega[\gamma\delta],
\end{aligned}
\]

with one common sign. Multiplicativity gives the same integer in all
three formulas. When \(p=-1\), these maps are units. The spectral
sequence then yields \(H^1(X)=H^2(X)=H^3(X)=0\); universal coefficients
and Poincare duality give the integral homology of \(S^6\). Together
with \(\pi_1(X)=1\) and \(\Theta_6=0\), this implies
\(X\cong_{\mathrm{diff}}S^6\).

This route does not use the global collapse theorem. Its
theorem-critical unresolved inputs are instead: the based meridian
relations, the integral specialization theorem at \(W\), the
finite-fibre lattices, and the geometric realization of the first
transgression and its multiplicative comparison.

\subsubsection{The direct algebraic-dimension
chain}\label{the-direct-algebraic-dimension-chain}

For an integral class

\[
E=E_{\gamma u}\gamma u+E_{\gamma w}\gamma w+E_{\gamma\delta}\gamma\delta
  +E_{uw}uw+E_{u\delta}u\delta+E_{w\delta}w\delta,
\]

Section 9.1 proves

\[
E\in\operatorname{NS}(F_b)
\quad\Longleftrightarrow\quad
P_E(z)=0,
\]

where

\[
P_E=E_{\gamma u}-E_{\gamma w}\tau-E_{\gamma\delta}\mu
 +6E_{uw}\mu+E_{u\delta}\beta+E_{w\delta}(6\mu^2-\tau\beta).
\]

The transformation laws under the two orbifold generators prove that

\[
1,\ \tau,\ \mu,\ \beta,\ 6\mu^2-\tau\beta
\]

are linearly independent. Hence the very general fibre has

\[
\operatorname{NS}(F_b)=\mathbb Z\eta,
\qquad
\eta=uw+6\gamma\delta.
\]

Moreover

\[
\eta^2=12,
\qquad
G_\eta=
\begin{pmatrix}
2\operatorname{Im}\tau&2\operatorname{Im}\mu\\
2\operatorname{Im}\mu&\frac13\operatorname{Im}\beta
\end{pmatrix},
\]

and

\[
\det G_\eta=\frac{2}{3}\operatorname{Im}\tau
\left(\operatorname{Im}\beta-\frac{6(\operatorname{Im}\mu)^2}{\operatorname{Im}\tau}\right)<0.
\]

Thus every nonzero multiple of \(\eta\) is nondegenerate and indefinite,
while its square is nonzero. A two-torus of algebraic dimension two
would have an ample positive class; one of algebraic dimension one would
have a nonzero fibre class of square zero. Neither can occur when
\(\operatorname{NS}=\mathbb Z\eta\). Therefore the very general fibre
has algebraic dimension zero.

A separate descent lemma shows that a meromorphic function on \(X\)
which is constant on uncountably many smooth fibres comes from the base.
Since very general fibres have algebraic dimension zero,

\[
\mathcal M(X)=f^*\mathcal M(\mathbb P^1)=f^*\mathbb C(t),
\qquad
a(X)=1,
\]

and \(f\) is the algebraic reduction. This proof is independent of the
topology and of the CDP dispute.

\subsubsection{What Section 10 has and has not
established}\label{what-section-10-has-and-has-not-established}

The local conductor calculation produces a nonzero ambient cotangent
section on a reduced nonnormal normal-crossing fibre even when the
torsion-free intrinsic cotangent sheaf has no section. For a line bundle
\(A\) trivial after normalization, the pair

\[
\bigl(\eta^*(df|_W)\otimes e,0\bigr)
\]

descends through the normalization/conductor sequence to

\[
s\in H^0(W,\Omega_X^1|_W\otimes A),
\]

whose intrinsic image is conductor-supported torsion. This local
mechanism is mathematically real and was independently reconstructed.

Its application to the global contradiction requires the exact analytic
chain

\[
H^2(W,E)^\vee
\simeq \operatorname{Hom}_W(E,\omega_W)
\simeq H^0(W,E^\vee\otimes\omega_W),
\qquad E=(T_X\otimes L)|_W,
\]

followed by the top-degree base-change map

\[
(R^2f_*(T_X\otimes L))_0\otimes k(0)
\longrightarrow H^2(W,(T_X\otimes L)|_W).
\]

The manuscript states the relevant Gorenstein duality and top-degree
base-change hypotheses, and no internal contradiction was found.
Nevertheless, three assertions must remain distinct:

\begin{enumerate}
\def\labelenumi{\arabic{enumi}.}
\tightlist
\item
  the local computation may expose a gap in a printed proof at a
  nonnormal fibre;
\item
  that does not by itself prove the published theorem false in all
  stated generality;
\item
  the theorem is contradicted by this example only after the
  construction and every hypothesis on \(X\) have independently been
  certified.
\end{enumerate}

\subsection{Revised status ledger}\label{revised-status-ledger}

\begin{longtable}[]{@{}
  >{\raggedright\arraybackslash}p{(\columnwidth - 4\tabcolsep) * \real{0.3333}}
  >{\raggedright\arraybackslash}p{(\columnwidth - 4\tabcolsep) * \real{0.3333}}
  >{\raggedright\arraybackslash}p{(\columnwidth - 4\tabcolsep) * \real{0.3333}}@{}}
\toprule\noalign{}
\begin{minipage}[b]{\linewidth}\raggedright
Status
\end{minipage} & \begin{minipage}[b]{\linewidth}\raggedright
Claim
\end{minipage} & \begin{minipage}[b]{\linewidth}\raggedright
Present assessment
\end{minipage} \\
\midrule\noalign{}
\endhead
\bottomrule\noalign{}
\endlastfoot
\textbf{VERIFIED exact algebra} & monodromy, Smith forms,
invariant/coinvariant lattices, period transformations, \(P_E\),
\(\eta^2\), Gram determinant, Appendix B rank matrices & Reproduced in
two local harnesses: 28 legacy checks and 14 supplemental checks, all
passing. \\
\textbf{CONDITIONALLY VALIDATED} & direct proof of \(a(X)=1\) and
\(\mathcal M(X)=f^*\mathbb C(t)\) & The complete argument is present and
survives line-by-line review, conditional on the period-family
construction and standard complex-torus structure theory. \\
\textbf{CONDITIONALLY VALIDATED} & Leray route from specialization and
transgression data to \(H_*(X)=H_*(S^6)\) & Formal sheaf,
spectral-sequence, rank, saturation, and multiplicativity logic survives
review. Geometric bridge inputs remain open to specialist
replication. \\
\textbf{OPEN} & existence and exact gluing of the global compact complex
threefold & No fatal defect found, but the infinite toric quotient,
finite logarithmic fillings, and Hausdorff/proper gluing remain
construction-level obligations. \\
\textbf{OPEN} & Section 10 as a counterexample to the printed CDP
conclusions & Local conductor mechanism is strong; global duality/base
change and exact comparison with the published proof require specialist
adjudication. \\
\textbf{REFUTED local wording} & Remark 7.28 as a fixed-triple
biconditional & False; must be existential in \(\ell_0\). \\
\textbf{REFUTED proof sentence} & Proposition 4.7(iv)'s properness
justification & Closedness needs the closed-graph/section argument. \\
\textbf{REFUTED notation/variance} & Lemma A.3's displayed
Mayer--Vietoris sequence as homology & It is cohomology. The stated
groups survive after correction and UCT. \\
\end{longtable}

\section{Disposition}\label{disposition}

This addendum is incorporated verbatim into the primary revised report.
The full section-by-section audit, issue ledger, proof-obligation
analysis, and references follow there.

\end{document}
